\documentclass[conference]{IEEEtran}
\IEEEoverridecommandlockouts

\renewcommand{\baselinestretch}{0.87}
% The preceding line is only needed to identify funding in the first footnote. If that is unneeded, please comment it out.
\usepackage{cite}
\usepackage{amsmath,amssymb,amsfonts}


\usepackage{float}
\usepackage{placeins}
\usepackage{algorithmic}
\usepackage{graphicx}
\usepackage{comment}
\usepackage{textcomp}
\usepackage{xcolor}
\usepackage[pdftex]{hyperref}
\usepackage{hyperref}
\hypersetup{hidelinks}
\usepackage{makecell}
\usepackage{multirow}
\usepackage{dingbat}
\usepackage{float}
\usepackage{biblatex}
\addbibresource{bibliography.bib}

\setlength{\arrayrulewidth}{0.2mm}
\setlength{\tabcolsep}{2pt}
\renewcommand{\arraystretch}{1.5}

\usepackage[margin=15mm]{geometry}
\usepackage{tabularray}

\usepackage{makecell, longtable}
\renewcommand\theadfont{\bfseries\small}
\renewcommand\theadgape{}
\usepackage{stfloats}
    \renewcommand{\dblfloatpagefraction}{.9}

\usepackage{amsmath}
\usepackage[skip=1ex]{caption}

\usepackage{lipsum}


\def\BibTeX{{\rm B\kern-.05em{\sc i\kern-.025em b}\kern-.08em
    T\kern-.1667em\lower.7ex\hbox{E}\kern-.125emX}}

    
\begin{document}

\title{Interpretable by Design, Efficient by Distillation: An Industrial Study of Expert-Guided Search-Based Software Team Formation [Supplementary Material]\\
}


\author{
\IEEEauthorblockN{
Felipe Cunha\IEEEauthorrefmark{1},
Mirko Perkusich\IEEEauthorrefmark{1},
Rian Gabriel Santos Pinheiro\IEEEauthorrefmark{2}\\
Evandro Costa\IEEEauthorrefmark{2},
Danyllo Albuquerque\IEEEauthorrefmark{1},
Kyller Gorgônio\IEEEauthorrefmark{1},
Angelo Perkusich\IEEEauthorrefmark{1}
}

\IEEEauthorblockA{
\IEEEauthorrefmark{1}VIRTUS Research Group, Federal University of Campina Grande, Campina Grande, Brazil\\
\IEEEauthorrefmark{2}Federal University of Alagoas, Maceió, Brazil
}
}

\maketitle

\renewcommand{\thetable}{S\arabic{table}}

\section{Introduction}

This document provides supplementary material for the paper \emph{Interpretable by Design, Efficient by Distillation: An Industrial Study of Expert-Guided Search-Based Software Team Formation}. It preserves detailed experimental evidence summarized in the main paper without altering or recalculating any data.

Table~\ref{tab:data-summary} characterizes the industrial dataset used to calibrate and evaluate the pipeline, including the candidate pool, technical taxonomy, temporal project split, team-size distribution, and collaboration graph. Table~\ref{tab:rq5-design} documents the retained practitioner-validation design and distinguishes unique cases, repeated comparisons used to examine consistency, and isolated acceptability judgements.

Table~\ref{tab:rq2-ilp-ga-summary} provides the complete project-level evidence for RQ2 by comparing the exact MILP reference and the calibrated GA through solution quality, stochastic variability, optimum-hit counts, runtime, and speedup. Table~\ref{tab:rq3-backtest} provides the complete project-level evidence for RQ3 by comparing the calibrated-GA recommendations with the historical teams through overall team fit, collaborative adequacy, and relative gain on the sixteen held-out 2025 projects.

Together, these tables support reproducibility and allow readers to trace the aggregate results and interpretations reported in the paper to their underlying project-level evidence.

%aqui estão as tabelas do artigo do SEIP, não mexa nos dados:


\begin{table}[H]
\centering
\footnotesize
\setlength{\tabcolsep}{3pt}
\caption{Summary of the industrial dataset.}
\label{tab:data-summary}
\begin{tabular}{p{0.34\linewidth}p{0.58\linewidth}}
\hline
Aspect & Summary \\
\hline
Candidate pool & 512 developers with at least one classified technical competence. \\
Technical taxonomy & 13 domains, 213 ecosystems, and 23 programming languages.
Median per developer: 2 domains, 6 ecosystems, and 3 languages. \\
Calibration projects & 205 projects starting up to 2024. Team sizes:
$k{=}4$: 45, $k{=}5$: 38, $k{=}6$: 37, $k{=}7$: 30, $k{=}8$: 26, $k{=}9$: 29. \\
Held-out projects & 16 projects starting in 2025. Team sizes:
$k{=}4$: 4, $k{=}5$: 4, $k{=}6$: 1, $k{=}7$: 1, $k{=}8$: 2, $k{=}9$: 4. \\
Collaboration graph & Pairwise collaboration history extracted from prior project co-participation and transformed through the expert-calibrated compatibility function. \\
\hline
\end{tabular}
\end{table}


\begin{table}[H]
\caption{Retained practitioner-validation design.}
\label{tab:rq5-design}
\centering
\small
\resizebox{\columnwidth}{!}{%
\begin{tabular}{lrrr}
\hline
Task & Unique cases & Repeats & Rows \\
\hline
Calibrated GA vs. historical & 16 & 3 & 19 \\
Calibrated GA vs. uniform GA & 12 & 2 & 14 \\
Calibrated-GA acceptability & 16 & -- & 16 \\
\hline
\end{tabular}
}
\end{table}

\clearpage
\begin{table*}[!t]
\centering
\scriptsize
\setlength{\tabcolsep}{4pt}
\caption{RQ2 comparison between the exact MILP reference and the calibrated GA on the twelve size-four controlled benchmark projects.}
\label{tab:rq2-ilp-ga-summary}
\begin{tabular}{lrrrrrrr}
\hline
Project & MILP AE & GA best AE & Best gap (\%) & GA AE mean $\pm$ SD & Hits/30 & MILP time (s) & Speedup \\
\hline
P1  & 0.818783 & 0.818783 & 0.00 & 0.788474 $\pm$ 0.050642 & 12 & 109.69  & 7.85$\times$ \\
P2  & 0.844445 & 0.844445 & 0.00 & 0.769284 $\pm$ 0.059689 & 7  & 208.44  & 10.26$\times$ \\
P3  & 0.813492 & 0.810847 & 0.33 & 0.789697 $\pm$ 0.032128 & 0  & 186.75  & 7.61$\times$ \\
P4  & 0.864286 & 0.864286 & 0.00 & 0.817566 $\pm$ 0.033847 & 1  & 45.74   & 4.41$\times$ \\
P5  & 0.843915 & 0.826720 & 2.04 & 0.771497 $\pm$ 0.032369 & 0  & 34.46   & 2.20$\times$ \\
P6  & 0.781746 & 0.781746 & 0.00 & 0.753527 $\pm$ 0.035620 & 14 & 508.73  & 93.51$\times$ \\
P7  & 0.870900 & 0.859326 & 1.33 & 0.858016 $\pm$ 0.004189 & 0  & 585.89  & 117.37$\times$ \\
P8  & 0.857672 & 0.857672 & 0.00 & 0.790728 $\pm$ 0.054681 & 10 & 195.47  & 27.33$\times$ \\
P9  & 0.867725 & 0.867725 & 0.00 & 0.854559 $\pm$ 0.023670 & 16 & 122.19  & 15.06$\times$ \\
P10 & 0.900794 & 0.900794 & 0.00 & 0.876323 $\pm$ 0.024483 & 11 & 105.33  & 26.60$\times$ \\
P11 & 0.834656 & 0.834656 & 0.00 & 0.779374 $\pm$ 0.031629 & 2  & 73.92   & 12.84$\times$ \\
P12 & 0.877513 & 0.857672 & 2.26 & 0.794899 $\pm$ 0.028884 & 0  & 226.69  & 45.79$\times$ \\
\hline
Summary & -- & -- & 0.50 mean & -- & 73/360 & -- & 13.95$\times$ median \\
\hline
\end{tabular}
\end{table*}


\begin{table*}[!t]
\centering
\footnotesize
\setlength{\tabcolsep}{4pt}
\caption{RQ3 backtest on the sixteen held-out 2025 projects.}
\label{tab:rq3-backtest}
\begin{tabular}{lrrrrrr}
\hline
Project & $k$ & $\mathrm{AE}_{\mathrm{real}}$ & $\mathrm{AE}_{\mathrm{GA}}$ & Gain (\%) & $\mathrm{AC}_{\mathrm{real}}$ & $\mathrm{AC}_{\mathrm{GA}}$ \\
\hline
P1  & 9 & 0.5238 & 0.6114 & 16.72 & 0.5143 & 0.8111 \\
P2  & 8 & 0.3750 & 0.5167 & 37.78 & 0.3643 & 0.6000 \\
P3  & 5 & 0.4532 & 0.6822 & 50.53 & 0.4200 & 0.7200 \\
P4  & 9 & 0.4008 & 0.6506 & 62.32 & 0.3571 & 0.8167 \\
P5  & 4 & 0.5762 & 0.6365 & 10.47 & 0.6000 & 0.8667 \\
P6  & 7 & 0.5246 & 0.5393 & 2.81  & 0.7667 & 0.7952 \\
P7  & 5 & 0.5817 & 0.6326 & 8.75  & 0.6333 & 0.8600 \\
P8  & 5 & 0.4365 & 0.6757 & 54.80 & 0.4000 & 0.8400 \\
P9  & 6 & 0.4048 & 0.5311 & 31.20 & 0.4000 & 0.6467 \\
P10 & 8 & 0.5869 & 0.6162 & 5.00  & 0.6643 & 0.8429 \\
P11 & 9 & 0.3274 & 0.5041 & 53.98 & 0.3071 & 0.6833 \\
P12 & 4 & 0.3770 & 0.5178 & 37.35 & 0.3667 & 0.7667 \\
P13 & 4 & 0.7817 & 0.8643 & 10.56 & 0.7667 & 0.9000 \\
P14 & 4 & 0.4008 & 0.9008 & 124.75 & 0.3000 & 0.9000 \\
P15 & 5 & 0.3770 & 0.5347 & 41.84 & 0.3667 & 0.7400 \\
P16 & 9 & 0.4373 & 0.4963 & 13.48 & 0.4810 & 0.7667 \\
\hline
Mean & -- & 0.4728 & 0.6194 & 35.15 & 0.4818 & 0.7847 \\
\hline
\end{tabular}
\end{table*}

\FloatBarrier
\clearpage

\end{document}
